% BHU PYQ -- Professional Exam Template
\documentclass[11pt,a4paper]{article}

\usepackage[a4paper,top=2.5cm,bottom=2.5cm,left=2.8cm,right=2.8cm,headheight=14pt]{geometry}
\usepackage{amsmath,amssymb}
\usepackage[T1]{fontenc}
\usepackage[utf8]{inputenc}
\usepackage{lmodern}
\usepackage{microtype}
\usepackage{fancyhdr}
\usepackage{enumitem}
\usepackage{listings}
\usepackage{hyperref}
\usepackage{lastpage}
\usepackage{parskip}

\hypersetup{colorlinks=true,urlcolor=black,linkcolor=black,
  pdftitle={CS-104: Computer Organization \& Architecture -- BHU PYQ}}

\pagestyle{fancy}
\fancyhf{}
\renewcommand{\headrulewidth}{0.4pt}
\renewcommand{\footrulewidth}{0.4pt}
\fancyhead[L]{\small   \textit{Banaras Hindu University}}
\fancyhead[R]{\small   \textit{CS-104: Computer Organization \& Architecture}}
\fancyfoot[C]{\small Page  \thepage\ of \pageref{LastPage}}

\lstset{
  basicstyle=\small tfamily,
  frame=single,
  breaklines=true,
  columns=flexible,
  keepspaces=true
}


\newlist{parts}{enumerate}{2}
\setlist[parts,1]{label=(\alph*),leftmargin=*,itemsep=4pt,topsep=3pt}
\setlist[parts,2]{label=(\roman*),leftmargin=*,itemsep=2pt,topsep=2pt}

\newcommand{\pts}[1]{\hfill   \textit{[#1]}}


\begin{document}

\begin{center}
  {\large\bfseries BANARAS HINDU UNIVERSITY}\\[2pt]
  {\normalsize B.Sc.\ (Hons.) Semester IV Examination, 2022--23}\\[6pt]
  \rule{0.8\linewidth}{0.5pt}\\[6pt]
  {\large\bfseries Paper: CS-104 --- Computer Organization \& Architecture}\\[4pt]
  {\normalsize Time: Three Hours \hfill Maximum Marks: 70}
\end{center}

\rule{\linewidth}{0.4pt}

\smallskip



\noindent   \textit{Instructions: Attempt any   \textbf{five} questions including Question~1 which is compulsory.}

\bigskip




\noindent   \textbf{Question 1.}\pts{$5+5=10$}

\begin{parts}
  \item State True or False and justify briefly:
  
\begin{enumerate}[label=(\roman*)]
    \item All high-level languages are compiled.
    \item All low-level languages are compiled.
    \item Storage is not a part of data processing.
    \item Concept of jitter is related to audio-video data.
    \item Port address is used for source to destination delivery.
  \end{enumerate}
  \item Briefly define the following:
  
\begin{enumerate}[label=(\roman*)]
    \item Malware
    \item Browser Hijacker
    \item Macro Virus
    \item FTP
  \end{enumerate}
\end{parts}

\bigskip




\noindent   \textbf{Question 2.}\pts{$10+5=15$}

\begin{parts}
  \item What is data processing? What are its different types? Explain each of them briefly.
  \item Explain the concept of stored program computers. Why is it considered better than hard-wired computers?
\end{parts}

\bigskip




\noindent   \textbf{Question 3.}\pts{$10+5=15$}

\begin{parts}
  \item What do you understand by a web crawler? Explain the working of a web crawler. What type of pages cannot be accessed by a web crawler?
  \item What is personalization of search results? State some of the parameters used for personalization of search results.
\end{parts}

\bigskip




\noindent   \textbf{Question 4.}\pts{$10+5=15$}

\begin{parts}
  \item Explain the different types of network topologies with examples. Also give their advantages and disadvantages.
  \item Explain the different components of a communication system.
\end{parts}

\bigskip




\noindent   \textbf{Question 5.}\pts{$10+5=15$}

\begin{parts}
  \item What is the role of Internet Service Providers (ISPs) in the Internet? Explain their different types.
  \item Differentiate between protocol and standards. What are the different types of standards used in communication? Name any one standard organization.
\end{parts}

\bigskip




\noindent   \textbf{Question 6.}\pts{$10+5=15$}

\begin{parts}
  \item Differentiate between a client and a server. Explain the roles of an e-mail client and e-mail server. How does message transfer take place in an e-mail communication system? Explain.
  \item Explain the applications of ICT in Education.
\end{parts}

\bigskip




\noindent   \textbf{Question 7.}\pts{$10+5=15$}

\begin{parts}
  \item Explain the memory hierarchy of a computer system.
  \item Differentiate between a CPU and a GPU. List some tasks where a GPU performs better than a CPU and explain why.
\end{parts}

\vfill
\rule{\linewidth}{0.4pt}

\begin{center}\small
\href{https://bhu-exam.vercel.app/}{Click here to visit our website}\\[4pt]\href{https://me-aryan.vercel.app/}{Click here to visit our contributor}\\[4pt]\href{https://quashan-me.vercel.app/}{Click here to visit our contributor}
\end{center}
\end{document}